\section{Scope: from one sextic triangle to arbitrary mixed components}
Turn 3 supplied an actual closed canonical triangle with constant primitive sextic
quotient, and separately proved square-source escape for all-$3\bmod4$ triangles
\cite{Turn3}. Neither canonical cycle closure nor cubic reciprocity forces an ES
witness. The present continuation retains the square-source mechanism but removes
its stabilizer-index and vertex-type restrictions. It proves an explicit height
bound for a square-source-closed set of any prescribed cardinality, a finite
cofactor reduction, and a computer-assisted exclusion of all such sets of at most
five vertices. The proof is independent of whether the local selectors are empty.
It does not prove that every component contains an occupied selector.

The canonical factor graph is inherited from Bryan \cite[\S\S1--7]{Bryan}.
Quadratic reciprocity is the classical law in Hatcher \cite[\S6.4, pp.205--212]{Hatcher}.
All counting and finite reduction arguments used below have complete proofs here.
A different Type-A/B square-lift shadow theorem \cite[\S\S3--6]{Shadow} allows
unbounded auxiliary parameters; it does not supply the fixed-$p$, source-below-$p$
conclusions proved here. No historical-priority conclusion is asserted.

\section{Every original square source and its coefficient return}
Let $p$ be prime in $1,121,169,289,361,529\pmod{840}$ and let
\[
 V_p=\{q<p:q\text{ prime},\ (q/p)=-1\},\quad
 \sigma_q=\begin{cases}1&q\equiv3\pmod4,\\3&q\equiv1\pmod4.\end{cases}
\]
Here $p\ge1009$, $V_p\ne\varnothing$, and its primes are at least 11.
Indeed $2,3,5,7$ are residues at each hard class. Factoring a least positive
nonresidue gives a nonresidue prime below $p$. For an odd integer $t>0$ retain
\[
 R_0=\sigma_q q,\quad R=R_0t^2<3p,\quad A_t=(p+R)/4.
                                                        \tag{1}
\]
Then $p/4<A_t<p$, $\gcd(pA_t,R)=1$, and $(A_t/p)=-1$. The last equality follows
from $4A_t\equiv\sigma_q q t^2\pmod p$; the bound implies $t<p$, and
$(\sigma_q/p)=1$. Thus the complete factorization of $A_t$ has at least one
nonresidue prime $r<p$, distinct from $q$. Every such prime is an outgoing edge,
with its actual weight $v_r(A_t)$; the sum of those weights is odd.

At each separately marked source retain
\[
 A_t=\prod_l l^{e_l},\quad \beta_l\in[-e_l,e_l],\quad
 u=\prod_l l^{e_l+\beta_l}\mid A_t^2,\quad
 v(\beta)=uA_t^{-1}\pmod R.                             \tag{2}
\]
The original gates are
\[
 E:R\mid4u+1,\qquad M:R\mid4u+p\iff R\mid u+A_t.       \tag{3}
\]
The map (2) is bijective, with $\beta_l=v_l(u)-e_l$. Its integer coefficients count
these original vectors, not a multiset of independently coloured prime copies.

For a hit put $d=\gcd(A_t,u)$ and
\[
 (h,r,s)=(d^2/u,u/d,A_t/d),\quad A_t=hrs,\quad u=hr^2.
\]
Prime valuations show $h,r,s>0$ and $\gcd(r,s)=1$. The ordered returns are
\[
 E:(x,y,z)=(A_t,hs\kappa,phr\kappa),\quad\kappa=(pr+s)/R,
\]
\[
 M:(x,y,z)=(A_t,phs\lambda,phr\lambda),\quad\lambda=(r+s)/R. \tag{4}
\]
The gates imply integrality because $h,r,s,p$ are units modulo $R$.
Multiplying the reciprocal sums by $phrs\kappa$ or $phrs\lambda$ proves $4/p$.
The ordered divisor inverses are $u=A_t^2/(Ry-pA_t)$ for E and
$u=pA_t^2/(Ry-pA_t)$ for M. These are inherited Type I/II coordinates
\cite[\S2, Propositions 2.2 and 2.6]{ET}, rechecked here with (1).
An upper-half source cannot have an M hit, since its other two denominators
are multiples of $p$. An E hit with $A_t>p/2$ returns to the first-half chart
by $(h,r,s,\kappa)\mapsto(h,\kappa,s,r)$, retaining the swap of the first two
denominators \cite{Turn3}.
The literal Elsholtz--Tao coordinate maps are
$(r,\kappa,s,h,(r+\kappa)/s,R)$ for E, with their image ordered $(z,x,y)$,
and $(r,s,\lambda,h,R,4rh\lambda-1)$ for M, with image $(x,z,y)$.
The inverse reads these coordinates and applies the displayed permutation;
$p(r+\kappa)=s(4hr\kappa-1)$ proves the extra E integrality.

\begin{theorem}[Exact carry for a square residual]\label{thm:carry}
For $c\in\{E,M\}$ let $\eta_E=1$, $\eta_M=p$. On the original base-candidate set
\[
 \mathcal B_{c,q,t}=\{u\mid A_t^2:R_0\mid4u+\eta_c\}
\]
put
\[
 k_c(u)=\frac{4u+\eta_c}{R_0},\quad d_c(u)=k_c(u)\bmod t^2,\quad
 b_c(u)=\frac{k_c(u)-d_c(u)}{t^2}.
\]
The full channel count is the exact integer coefficient
\[
 \boxed{|c_{q,t}|=[Z^0]\sum_{u\in\mathcal B_{c,q,t}}Z^{d_c(u)}}. \tag{5}
\]
The inverse, with the source and channel retained, is
\[
 \boxed{u=\frac{R_0(d_c+t^2b_c)-\eta_c}{4}.}             \tag{6}
\]
It applies exactly when the right side is a positive integer divisor of $A_t^2$.
No coprimality of $R_0$ and $t$ is required.
\end{theorem}
\begin{proof}
Euclidean division gives $k=d+t^2b$ uniquely. Divisibility by $R_0t^2$ is exactly
$d=0$, proving (5); solving the same integer equation gives (6). If only $d$ is
kept, its full fibre is the displayed set of allowed $b$ values. Keeping a separate
basis vector for each original $u$ makes that fibre explicit.
\end{proof}

The coefficient pushforward $\mathbb Z^{\mathcal B}\to\mathbb Z^{\mathbb Z/t^2}$
sums along the actual fibres. Its kernel is generated by $e_u-e_{u_*}$ in each
occupied fibre, where $u_*$ is its least divisor. Recording all nonselected
coordinates and the sum reconstructs the selected coordinate by subtraction.
A coefficient that becomes zero after reduction remains a supported zero, whereas
an empty original fibre is absent. No inverse of the augmentation is asserted.

For $t=2j+1$ write $F_q(j)=A_t$. Then
\[
 F_q(j)=A_1+\sigma_q q\,j(j+1),\qquad
 \boxed{\gcd(F_q(i),F_q(j))=\gcd(F_q(i),(j-i)(i+j+1)).}  \tag{7}
\]
This follows by subtraction and $\gcd(F_q(i),\sigma_q q)=1$.
The exact unchanged-divisor domain of two sources is
$\operatorname{Div}(\gcd(A_t,A_{t'})^2)$.
On it the centred coordinates change by
$\beta_l' =\beta_l+v_l(A_t)-v_l(A_{t'})$, and the integer carry $k_c(u)$ is
unchanged. The full tests remain $t^2\mid k_c(u)$ and $t'^2\mid k_c(u)$.
This is a coefficient transport through a common original source; it does not
assert that a nonzero projected coefficient has a full lift.

\section{Paired roots and a uniform source that avoids a finite vertex set}
For $r\in V_p\setminus\{q\}$ the discriminant of $F_q$ modulo $r$ is
$-\sigma_q q p\ne0$. Consequently
\[
 \nu_{q,r}=\#\{j\pmod r:F_q(j)=0\}
 =1+\left(\frac{-\sigma_q q p}{r}\right)\in\{0,2\}.    \tag{8}
\]
When roots exist they are exchanged by $j\mapsto-1-j$, with no fixed root.
At $r=q$ there is no root. These are statements about prime divisibility of the
actual integer source, not about availability of a selector divisor.

For $C\subseteq V_p$ containing $q$, define
\[
 Z_C(q,N)=\#\{0\le j<N:\gcd(F_q(j),\prod_{r\in C}r)=1\}.
\]
Let $\mathcal R$ contain the $n$ primes of $C\setminus\{q\}$ for which (8) is two,
and put $\delta=\prod_{r\in\mathcal R}(1-2/r)$. For $J\subseteq\mathcal R$,
let $d_J=\prod_{r\in J}r$ and let $\mathcal Z_J\subset[0,d_J)$ be the complete
CRT root set. The empty set has modulus one and root zero. Inclusion--exclusion is
\[
 \boxed{Z_C(q,N)=\sum_{J\subseteq\mathcal R}(-1)^{|J|}
 \sum_{a\in\mathcal Z_J}\left(1+\left\lfloor\frac{N-1-a}{d_J}\right\rfloor\right).} \tag{9}
\]
Each term is a count of original source indices.

\begin{theorem}[Reflection bound and explicit escape prefix]\label{thm:prefix}
If $n>0$ then
\[
 \boxed{|Z_C(q,N)-N\delta|<(3^n-1)/2,\qquad
 \delta\ge\frac9{2n+9}.}                              \tag{10}
\]
For $n=0$, $Z_C(q,N)=N$. Define
\[
 (b_0,\ldots,b_7)=(1,2,3,4,5,6,9,15),\quad
 b_n=\left\lceil\frac{(2n+9)(3^n-1)}{18}\right\rceil\quad(n\ge8).
                                                               \tag{11}
\]
Then $Z_C(q,b_n)>0$. In particular, if $|C|=m$ and
\[
 \sigma_q q(2b_{m-1}-1)^2<3p,
\]
one of those original sources has a nonresidue prime factor outside $C$.
\end{theorem}
\begin{proof}
For nonempty $J$, its $2^{|J|}$ CRT roots are paired by $a\mapsto d_J-1-a$.
There is no fixed point: it would imply $2a+1=0$ at every prime of $J$, contrary
to $p\ne0$ there. For one pair, remove complete periods from $N$. In a prefix of
length at most $(d_J-1)/2$ at most one root is counted; in a longer proper prefix
at least one is counted. In either case the discrepancy from $2N/d_J$ has
absolute value strictly less than one. Thus the $J$ discrepancy is less than
$2^{|J|-1}$. Summing gives
$\sum_{J\ne\varnothing}2^{|J|-1}=(3^n-1)/2$.
This improved estimate concerns prefixes beginning at zero, not arbitrary intervals.
If the active primes are increasing, $r_i\ge2i+9$, so the telescoping product is
$\delta\ge\prod_{i=1}^n(2i+7)/(2i+9)=9/(2n+9)$.

For $n\ge8$, (10) and (11) force $Z_C>0$. For $1\le n\le5$, a prime at least
11 can cover at most one index in $[0,n]$: two distinct root representatives
would sum to $r-1\ge10$, whereas their maximum sum is $2n-1\le9$.
For $n=6$ and $N=9$, only 11 and 13 can cover two indices, so the maximum covered
count is $6+2=8$. For $n=7$ and $N=15$, primes 11 and 13 cover at most three
indices, 17,19,23 at most two, and all others at most one. Hence at most
$7+2+2+1+1+1=14$ indices are covered. The constants increase with $n$ and satisfy
$b_n\le4^n$: for the formula use $(2n+9)/18<(4/3)^n$, starting at $n=1$ and
inducting; strict inequality before rounding gives the integer bound.
Finally all indicated sources are below $p$ under the stated strict range bound.
Their character is $-1$, so the C-free one has an actual nonresidue prime factor
below $p$ outside $C$.
\end{proof}

Original edge weights can be counted exactly as well. If $x$ is a root modulo
$r^a$, its unique lift is
\[
 x'=x+r^ah,\qquad h\equiv-\frac{F_q(x)}{r^a}
                 \bigl(\sigma_q q(2x+1)\bigr)^{-1}\pmod r.       \tag{12}
\]
The derivative is a unit at every root. Keeping both lifted roots at every layer,
\[
 \boxed{\sum_{j=0}^{N-1}v_r(F_q(j))=
 \sum_{\substack{a\ge1\\r^a\le F_q(N-1)}}
 \sum_{x\in\mathcal Z(r^a)}
 \left(1+\left\lfloor\frac{N-1-x}{r^a}\right\rfloor\right).}     \tag{13}
\]
This is the indicator expansion of each valuation; it is finite and exact.
Neither even multiplicities nor repeated visits to a prime are deleted.

\section{Mixed-type reciprocity, including opposite edges}
For $q,r\in V_p$, $q\ne r$, root eligibility is
\[
 \left(\frac qr\right)=-\left(\frac{-\sigma_q}{r}\right).
                                                               \tag{14}
\]
Reciprocity and its supplementary laws show that two orientations can be eligible
simultaneously exactly for the unordered type pairs
\[
 \boxed{(1,1),\ (5,5),\ (1,7),\ (5,7)\pmod{12}.}        \tag{15}
\]
For these pairs either both orientations or neither is eligible. For every other
type pair exactly one orientation is eligible before the finite source bound.
For example, for $q\equiv1\pmod4$ use
$(-3/r)=+1$ at $r\equiv1,7\pmod{12}$ and $-1$ at $r\equiv5,11\pmod{12}$;
for $q\equiv3\pmod4$ use $(-1/r)$. These signs with reciprocity give all (15).
Eligibility does not itself prove a bounded source occurs.

If $C$ contains $n_j$ vertices of each type $j\in\{1,5,7,11\}$, the number of
its possible directed edges is at most
\[
 U(C)=\binom m2+\binom{n_1}2+\binom{n_5}2+n_7(n_1+n_5),\quad m=|C|.
                                                               \tag{16}
\]
If $L$ initial square sources are valid at every vertex and $\min C>2L-2$,
closure would imply $mL\le U(C)$. Indeed (7) prevents a prime of $C$ from occurring
at two ports at the same vertex, so each vertex has at least $L$ distinct
successors. The old all-$3\bmod4$ theorem is its special case $U(C)=\binom m2$.
Using that smaller bound for arbitrary mixed vertices would be false.

\section{A uniform reduction of every higher-index closed component}
Call a nonempty set $C\subseteq V_p$ \emph{square-source closed} if every
nonresidue prime factor of every valid source (1) at every $q\in C$ lies in $C$.
No condition on the local gate counts is included in this definition. An all-failure
closed component of the full square-source graph is such a set.

\begin{theorem}[Cardinality--height bound and finite cofactor reduction]
\label{thm:height}
Let $C$ be square-source closed and $m=|C|$. Then $m\ge2$. Put
$N=b_{m-1}$ and $H=(2N-1)^2$. Every $q\in C$ satisfies
\[
 \boxed{\sigma_q qH\ge3p,\qquad q\ge p/H,\qquad
 p<(4H)^m<16^{m^2}.}                                   \tag{17}
\]
Moreover $C$ contains a simple canonical cycle of some length $2\le l\le m$.
Write its original equations as
\[
 p+\sigma_iq_i=4c_iq_{i+1},\qquad i\pmod l.             \tag{18}
\]
After cyclically rotating to a least cofactor $c_0=v$, every coefficient lies in
\[
 \boxed{1\le v<(H+3)/4,\qquad
 v\le c_i\le\left\lfloor\frac{vH-1}{4v-3}\right\rfloor.}       \tag{19}
\]
Thus each fixed cardinality has a completely explicit finite source-coefficient
search, valid for all hard primes and every stabilizer quotient.
\end{theorem}
\begin{proof}
Self edges are impossible and canonical sources are always valid, so finite
closure supplies a simple cycle of length at least two. Theorem~\ref{thm:prefix}
forces the first inequality of (17); otherwise a source has a nonresidue factor
outside $C$. Since $\sigma_q\le3$, $q\ge p/H$.

Let $D=4^l\prod c_i$, $S=\prod\sigma_i$. Composition of (18) gives
\[
 (D-S)q_i=pT_i,\quad T_i>0,\quad D>S.
\]
Primality and $q_i\ne p$ imply $p\mid D-S$. Also $c_i=A_{q_i,1}/q_{i+1}<H$,
so $p\le D-S<(4H)^l\le(4H)^m$. Since $b_{m-1}\le4^{m-1}$,
$4H<16^m$, proving the last bound.

For (19), let $Q$ be the largest vertex of this cycle. Its incoming edge gives
$4vQ\le p+\sigma q_{\rm prev}<p+3Q$, so $Q<p/(4v-3)$.
Each canonical source is therefore less than $vp/(4v-3)$. Since every successor
is at least $p/H$, $c_i<vH/(4v-3)$. The largest allowed integer is precisely the
right side of (19), and $c_i\ge v$ forces $4v-3<H$.
\end{proof}

Here is the full inverse of the coefficient word. For each cyclic starting point
$i$, iterate
\[
 D_0=1,\quad T_0=0,\qquad
 D_{j+1}=4c_{i+j}D_j,\quad T_{j+1}=\sigma_{i+j}T_j+D_j.
\]
Let $T_i$ denote its final numerator. Set
\[
 h=\gcd(T_0,\ldots,T_{l-1}),\quad
 \boxed{p=(D-S)/h,\qquad q_i=T_i/h.}                    \tag{20}
\]
For an original simple prime cycle this follows from $(D-S)q_i=pT_i$ and
$\gcd_iq_i=1$. Conversely (20) supplies at most one integer candidate; check
integrality, hard class, all primalities, distinctness, $11\le q_i<p$,
vertex types, all character values and every equation (18). Canonical cyclic
rotation is retained, with possible repeated least-cofactor roots explicitly
counted as duplicate profiles rather than duplicate components.

\section{No square-source-closed set of at most five vertices}
\begin{theorem}[Finite, all-prime exclusion]\label{thm:five}
For every hard prime $p$, every nonempty set of at most five vertices has some
valid original square source with $t\in\{1,3,5,7,9\}$ and a nonresidue prime factor outside that set.
In particular every all-failure closed component in the full square-source graph
has at least six vertices.
\end{theorem}
\begin{proof}
Suppose $|C|\le5$ and all its valid sources with $t\in\{1,3,5,7,9\}$ have every nonresidue factor in $C$. Five initial ports at a vertex have pairwise
source gcd prime factors at most eight by (7). Since every vertex is at least 11,
these ports require five distinct successors, impossible in $C\setminus\{q\}$.
Thus $81\sigma_q q\ge3p$ for each vertex. Apply the finite reduction (18)--(20)
with $H=81$ and $2\le l\le5$. Its complete domain is
\[
 \sigma_i\in\{1,3\},\quad 1\le v\le20,\quad c_0=v,\quad
 v\le c_i\le\left\lfloor\frac{81v-1}{4v-3}\right\rfloor.        \tag{21}
\]
Retain only profiles satisfying $81\sigma_iT_i\ge3(D-S)$ at every vertex.

The supplied standard-library finite proof enumerates (21), using the exact
monotone reduction described next. The complete count is 1,381,117,764 possible
rooted coefficient profiles, of which 5,922,259 satisfy all range inequalities.
There are 1,915 profiles reconstructing typed hard-class integer parameters.
Exact trial division leaves six all-prime profiles; the nonresidue tests leave
four rooted profiles, representing precisely the following three cycles:
\[
\begin{array}{r|l}
4201&61\to137\to1153\to383\to191\to61\\
12601&409\to3457\to5743\to2293\to487\to409\\
26209&661\to881\to7213\to5981\to5519\to661.
\end{array}                                                    \tag{22}
\]
All have five vertices, so $C$ must equal that cycle. The marked multiplier-nine
sources at their first vertices are respectively
\[
\begin{array}{r|r|l|r}
p&q&A_{q,3}&\text{outside nonresidue factor}\\\hline
4201&61&1462=2\cdot17\cdot43&43\\
12601&409&5911=23\cdot257&23\\
26209&661&11014=2\cdot5507&5507.
\end{array}                                                    \tag{23}
\]
Every source is below $p$. Every factor is prime and the last column has character
$-1$ modulo the corresponding $p$. These are checked by exhaustive trial division
and exact modular powers. Thus each candidate leaves its five vertices, contradicting
closure. The independently implemented enumerator recovers the identical complete
range-profile digest and candidate cycles. This is a finite computer-assisted
proof after (21), not extrapolation from a bounded prime census.
\end{proof}

\begin{corollary}[A terminating six-vertex arithmetic expansion]
At every hard prime, starting at any external nonresidue prime, factoring at most
25 valid original sources with $t\in\{1,3,5,7,9\}$ supplies at least six distinct
reachable external nonresidue primes. No factor inventory or target occupancy is
assumed.
\end{corollary}
\begin{proof}
Explore the five valid ports of each newly discovered vertex, retaining every
nonresidue factor. Stop when six distinct vertices have been discovered. If the
queue emptied sooner, the visited set would have size at most five and be closed
under those five ports, contradicting Theorem~\ref{thm:five}. Before six vertices
are found, at most five vertices, and hence at most 25 sources, can be processed.
Every factorization terminates by trial division; no polynomial running-time bound
is asserted. The earlier source labels and multiplicities remain attached to the
returned paths.
\end{proof}

The general height theorem has a similar constructive reading. Closure only under
the valid initial $b_{m-1}$ ports already suffices for its proof. Therefore, if
$p\ge16^{m^2}$, a breadth-first exploration of these ports from any starting vertex
must reach at least $m+1$ distinct vertices after at most $m b_{m-1}$ source
factorizations. This is proved finite factor growth, not termination at a selector.

\paragraph{Why the finite enumeration is complete.}
With other coefficients fixed and positive, let $c_j=x$. Then
\[
 D-S=D_0x-S,\qquad T_i=A_ix+B_i,\qquad A_i\ge0,\ B_i>0,
\]
so
\[
 \frac{d}{dx}\frac{T_i}{D-S}
 =\frac{-A_iS-D_0B_i}{(D_0x-S)^2}<0.                    \tag{24}
\]
The numerator expansion contains a positive term independent of every $c_j$,
which proves $B_i>0$. Thus the range-compatible coefficient set is downward
closed. At a partial profile, set all future coefficients to their least allowed
value $v$. If any inequality fails, the whole remaining box is excluded.
Otherwise the next coordinate is bounded by
\[
 [3D_0-81\sigma_iA_i]x\le81\sigma_iB_i+3S.              \tag{25}
\]
When the bracket is positive, its integer floor supplies the upper endpoint;
otherwise it imposes no upper restriction. Recursing through every integer up to
the least upper endpoint covers exactly the allowed profiles. One implementation
uses (25), the other tests (24) by binary search and reconstructs all $T_i$ from a
single expanded numerator and $4c_iT_{i+1}=D-S+\sigma_iT_i$.
The complete ledger keeps every one of the 1,915 typed models and a composite
factor or failed character whenever rejected. Its maximum reconstructed $p$ is
200,875,369. No primality assumption enters from a probable-prime test.

\section{An exact mixed-type local obstruction and its surviving path}
At the prime $p=12889\equiv289\pmod{840}$ the canonical component is
\[
 43\longrightarrow61\longrightarrow43,
\]
with $A_{43,1}=3233=53\cdot61$ and
$A_{61,1}=3268=2^2\cdot19\cdot43$. Each successor is the only nonresidue prime
factor and both E/M gates are empty. This is the allowed mixed pair of types
$7$ and $1\pmod{12}$, not a violation of reciprocity.

There is a stronger exact local counterexample: \emph{all} valid square sources
at each of these two vertices have both gates empty. At $43$, the full list is
$t=1,3,\ldots,29$; at $61$, it is $t=1,3,\ldots,13$. The certificate exhausts all
22 original divisor boxes and their carry coefficients. This disproves an
immediate-occupancy assertion even when every valid square at those vertices is
used; it does not contradict Theorem~\ref{thm:five}, which supplies outside factors.

A separately verified actual path is
\[
 43\xrightarrow{t=3}3319\xrightarrow{t=1}1013
 \xrightarrow{t=1}11,
\]
with source factorizations $3319$, $4052=2^2\cdot1013$,
$3982=2\cdot11\cdot181$, then $3225=3\cdot5^2\cdot43$.
The last source has an original M hit $u=9$, with
$(h,r,s,\lambda)=(1,3,1075,98)$, and returns
\[
 \frac4{12889}=\frac1{3225}+\frac1{1357856150}+\frac1{3789366}.
\]
The certificate retains the order and inverse. This is a specified successful path, not a theorem that
every outside factor eventually reaches a hit.

At $p=1129,q=11,t=11$, $R=1331$, $A_t=615=3\cdot5\cdot41$, and
$u=1845=3^2\cdot5\cdot41\mid A_t^2$. Its numerator is
\[
 4u+1=7381=121\cdot61.
\]
It satisfies the separate congruences modulo $11$ and $121$, but not their product
$1331$. The original carry is $671$, whose digit modulo $121$ is $66$.
Thus independent congruences impose only the lcm when the moduli overlap;
Theorem~\ref{thm:carry} retains the additional valuation exactly.

Finally C-free production is sufficient, not necessary, for an outside factor.
At $p=3361,q=977$, $A_{q,1}=1573=11^2\cdot13$ has both 11 and 13 as nonresidue
prime factors. Relative to $C=\{977,11\}$ it is not C-free, yet 13 is outside.
The exact function $Z_C$ is not substituted for the full closure predicate.

\section{Outcome and verification boundary}
The surviving square-source mechanism has a genuine uniform extension:
closed components of every prescribed cardinality reduce to an explicit finite
set of original integer cycle equations, at all stabilizer indices and both
vertex types. Every such component satisfies $p<16^{|C|^2}$, and all sets of
size at most five have a proved exit among their first five valid odd-square sources. These results do not bound $|C|$ above in terms
that contradict every possible hard prime. Large closed all-failure sets, and the
question whether the entire square-source family has a hit at each $p$, remain.

The original channel coefficient (5), not its augmentation or a coarse character,
is required for positivity. The quotient carry gives its complete reconstruction
including noncoprime square factors; it does not force that coefficient to be
nonzero. Turn 5 must combine the channel constraints on the remaining larger
sets or pursue the complete-shell sum, rather than equate source escape with an
occupied selector.

The complete finite cofactor search is the proof component of
Theorem~\ref{thm:five}. The separate bounded square-source scan covers the 13 hard
primes through 5000: 2698 vertices, 4460 valid sources, 125666 original divisors,
103 E states and 146 M states. It is corroboration, not an enlarged ES verification
range or evidence replacing the unbounded argument. The two code implementations
are independent implementations, not independent mathematical review. No Lean
build, historical-priority determination or resolution of ES is asserted.
